\section{Supporting delivery-mechanism probes}\label{app:delivery}
These probes connect documented tail-processor behavior to an explicit complete-witness predicate. A matched $2^4$ design crosses an early received root ERROR marker, an immediate or five-second-delayed decisive database leaf, decision caches enabled or disabled, and churn/capacity. Each combination repeats with three seeds and 200 seven-span failed traces. The decision wait is two seconds. Churn inserts 300 normal traces into 256 slots; its control has no churn and 50,000 slots. The churn factor therefore changes load and capacity jointly. Timing pairs share input content within other factors.

\begin{table}[htbp]
\centering\small
\begin{tabular}{llllrrr}
\toprule
Root ERROR & Leaf & Cache & Churn & IDs & Spans & Complete \\
\midrule
\tableinput{matched-probe-summary-rows.tex}
\bottomrule
\end{tabular}
\caption{Matched protected-only probes, per cell with 200 failed traces. Rows group equivalent factor outcomes; every reported result repeats in all three seeds. Complete witnesses require the stipulated span identities, parent links, and decisive leaf evidence. ``Either'' denotes both tested levels.}
\label{tab:matched}
\end{table}

An early received root ERROR preserves all complete witnesses in the tested combinations. Without that marker, immediate delivery preserves all witnesses while delayed delivery preserves none. With churn and no decision cache, delayed ERROR leaves can be exported without their discarded ancestors. A nonsampled cache can instead preserve the drop decision and export none. A separate burst probe with sixteen trace slots retains sixteen of 200 failures and records 184 early drops. The pinned processor documentation describes these decision, cache, and eviction mechanisms~\cite{otelTailSampling0136}.

Full configurations, additional probes, and partial-span interventions are retained in the artifact. Section~\ref{sec:discussion} states the delivery and evidence-predicate limitations.
